% One-time imported draft. Editable native source; no Markdown regeneration.
\section{Project contract}\label{003-eight-byte-sprite-assignment-draw-an-eight-byte-sprite-and-its-mirror}

Interpret eight source bytes as eight rows of an 8 by 8 monochrome
sprite. In each row, bit 7 is the leftmost pixel and bit 0 is the
rightmost. Draw the original in columns 0 through 7 of a 16 by 8
framebuffer. Reverse the bits in each source byte, store the result in
mirrored\_rows, and draw those reversed rows in columns 8 through 15. A
set bit must produce the RGB word 0x00ffffff; a clear bit must produce
zero.

The framebuffer begins at 0x10010000 in the default memory layout and
contains 128 words, or 512 bytes, in row-major order. Its row stride is
16 words, or 64 bytes. The supplied header configures a 128 by 64
display using 8 by 8 units. Preserve the source bytes and the initial
values of frame\_guard and rows\_guard. Produce no console output and
exit normally within the grader\textquotesingle s 10,000-instruction
budget per case.

Any of the eight source bytes may contain any pattern from 00 through
ff. The grader may replace the starter rows and poison the framebuffer
before execution. Write every required pixel, including black pixels.
The contract concerns final values and preservation; it does not promise
to detect temporary writes later restored before exit.

\subsection{Guided
  implementation}\label{003-eight-byte-sprite-assignment-guided-implementation}

Build the explicit eight-iteration reversal loop, then the column loop,
then the row loop. The guide uses named registers and checkpoint labels
so its traces can be compared. Each branch has an explicit nop delay
slot. Predict the state at the listed boundaries before stepping. The
instructor\textquotesingle s illustration checks concern this reference;
student grading does not require its scratch registers or private
labels.

\subsection{Free-writing
  implementation}\label{003-eight-byte-sprite-assignment-free-writing-implementation}

Use any equivalent implementation that preserves the published
interface. It must handle substituted rows, not merely reproduce the
starter image. The optional bitswap version is one example. Explain why
the implementation reverses each row\textquotesingle s eight bits and
why every framebuffer address remains within the allocated 512 bytes.

\subsection{Submission}\label{003-eight-byte-sprite-assignment-submission}

Submit main.asm, a reversal trace for 0x16, and a derivation of the row
stride and final pixel address. Explain the delay-slot behavior of a
branch used in the program. Run the public checks. A visible image is
useful evidence, but memory checks and a state trace can reveal errors
that a symmetric starter picture would conceal.
